\documentclass{article}
\usepackage{listings}

% if you need to pass options to natbib, use, e.g.:
%     \PassOptionsToPackage{numbers, compress}{natbib}
% before loading neurips_2026


% 6. "dblblindworkshop" option is used for the Workshop with double-blind reviewing
 \usepackage[dblblindworkshop]{neurips_2026}
% 7. "education" option is used for the Education Track
%  \usepackage[education]{neurips_2026}


% After being accepted, the authors should add "final" behind the track to compile a camera-ready version.
% 6. Workshop with double-blind reviewing
%  \usepackage[dblblindworkshop, final]{neurips_2026}
% "preprint" option is used for arXiv or other preprint submissions
 % \usepackage[preprint]{neurips_2026}

% to avoid loading the natbib package, add option nonatbib:
%    \usepackage[nonatbib]{neurips_2026}

\usepackage[utf8]{inputenc} % allow utf-8 input
\usepackage[T1]{fontenc}    % use 8-bit T1 fonts
\usepackage{hyperref}       % hyperlinks
\usepackage{url}            % simple URL typesetting
\usepackage{amsmath}        % for \lVert and \rVert
\usepackage{booktabs}       % professional-quality tables
\usepackage{amsfonts}       % blackboard math symbols
\usepackage{nicefrac}       % compact symbols for 1/2, etc.
\usepackage{microtype}      % microtypography
\usepackage{xcolor}         % colors
\usepackage{graphicx}

% Draft-only marker. Remove (or redefine to \gobble) before submission.
\newcommand{\todo}[1]{\textcolor{red}{[\textsc{todo}: #1]}}

% Notation shorthands.
\newcommand{\R}{\mathbb{R}}
\newcommand{\bx}{\mathbf{x}}
\newcommand{\bA}{\mathbf{A}}
\newcommand{\beps}{\boldsymbol{\varepsilon}}
\newcommand{\dototal}{\Delta_{\mathrm{total}}}
\newcommand{\dotraj}{\Delta_{\mathrm{traj}}}
\newcommand{\dooutcome}{\Delta_{\mathrm{outcome}}}
\newcommand{\Pa}{\mathrm{Pa}}
\newcommand{\dop}{\mathrm{do}}

% Note. For the workshop paper template, both \title{} and \workshoptitle{} are required, with the former indicating the paper title shown in the title and the latter indicating the workshop title displayed in the footnote. 
\workshoptitle{TAE (Trust-AI-Eval): Can We Trust AI Evaluation?}
\title{Does the World Agree? A Model-vs-World Audit for Counterfactual
Explanations of Time Series}


% The \author macro works with any number of authors. There are two commands
% used to separate the names and addresses of multiple authors: \And and \AND.
%
% Using \And between authors leaves it to LaTeX to determine where to break the
% lines. Using \AND forces a line break at that point. So, if LaTeX puts 3 of 4
% authors names on the first line, and the last on the second line, try using
% \AND instead of \And before the third author name.


\author{%
  David S.~Hippocampus\thanks{Use footnote for providing further information
    about author (webpage, alternative address)---\emph{not} for acknowledging
    funding agencies.} \\
  Department of Computer Science\\
  Cranberry-Lemon University\\
  Pittsburgh, PA 15213 \\
  \texttt{hippo@cs.cranberry-lemon.edu} \\
  % examples of more authors
  % \And
  % Coauthor \\
  % Affiliation \\
  % Address \\
  % \texttt{email} \\
  % \AND
  % Coauthor \\
  % Affiliation \\
  % Address \\
  % \texttt{email} \\
  % \And
  % Coauthor \\
  % Affiliation \\
  % Address \\
  % \texttt{email} \\
  % \And
  % Coauthor \\
  % Affiliation \\
  % Address \\
  % \texttt{email} \\
}


\begin{document}


\maketitle


\begin{abstract}
  The abstract paragraph should be indented \nicefrac{1}{2}~inch (3~picas) on both the left- and right-hand margins. Use 10~point type, with a vertical spacing (leading) of 11~points. The word \textbf{Abstract} must be centered, bold, and in point size 12. Two line spaces precede the abstract. The abstract must be limited to one paragraph.
\end{abstract}

\section{Introduction}

A counterfactual explanation answers the question practitioners actually ask of a high-stakes sequential model: \emph{what would have had to be different in this trajectory for the prediction to have been otherwise?} \citep{wachter2017counterfactual}. In the temporal setting the question is routinely posed about clinical trajectories, industrial telemetry and financial series, and a substantial family of methods now answers it \citep{delaney2021nativeguide,ates2021comte,hollig2022tsevo,li2024mcels,paredescetina2025confetti,yan2023counts}.

The evaluation protocol those methods are judged by has been remarkably stable: \emph{validity} (does the counterfactual flip the classifier), \emph{proximity} (how far is it from the factual), \emph{sparsity} (how many entries changed), and \emph{plausibility} (does it look like training data) \citep{pawelczyk2021carla,hollig2023xtscbench,nguyen2024amee}. Every one of those four quantities is a relation between the counterfactual and \textbf{the model}. None is a relation between the counterfactual and \textbf{the process that generated the data}. A method can therefore score perfectly on all four while proposing a trajectory the world could never have produced, or an intervention that would in fact have left the outcome untouched. The evaluation cannot tell the two apart, because it never asks the world.

This is a measurement-validity problem rather than a leaderboard problem, and it is the reason the gap has persisted. Asking whether the world agrees requires a setting where the world is known: a benchmark that ships the mechanism that generated its own data, so that the \emph{same} intervention a method claims credit for can be re-run through the true structural equations and the two outcomes compared. Real datasets cannot supply this, and existing synthetic temporal benchmarks were built for a different purpose, recovering the causal graph \citep{cheng2024causaltime,herdeanu2025causaldynamics,ferdous2025timegraph}, with no explanation method in the loop.

We build that setting and use it as an instrument. Our benchmark generates multivariate time series from a known temporal structural causal model (SCM) with additive noise, which makes Pearl's abduction step an exact subtraction \citep{pearl2009causality,hoyer2008anm}: the exogenous noise of a given instance is recovered rather than approximated, so the counterfactual the world would have produced \emph{for that instance} is computed exactly rather than estimated. On top of it we define a \textbf{model-vs-world audit}: a proposed counterfactual is admitted or rejected by a mechanism-consistency gate, its implied intervention is realised by the true mechanism, and the disagreement between what the explanation \emph{claims} and what the world \emph{does} is decomposed additively into a trajectory term (the proposed trajectory is not the one the world would produce) and an outcome term (the classifier does not track the true label rule).

Building the instrument surfaced a construct-validity issue that we consider part of the contribution rather than an implementation detail. The audit's answer depends on how a proposal is \emph{read}: interpreting a counterfactual as a single $\dop()$ at one timestep inflates the trajectory term for densely-editing methods, while interpreting every edited slice as its own intervention drives that term to zero \emph{by construction}. Neither endpoint is privileged, so we report the quantity that is stable across both; \textbf{do-complexity} $D$, the number of timesteps a proposal must declare as interventions before the mechanism can reproduce it, alongside every disagreement number. A method with $D$ close to the trajectory length has not explained a prediction; it has rewritten the trajectory, and its agreement with the world is vacuous because nothing was left for the mechanism to predict.

A second finding is a reporting rule rather than a method comparison. Stability of a temporal SCM requires a spectral radius below one(citation needed), so an intervention's effect on the label attenuates geometrically in the distance from the intervention time to the \emph{label site}, not to the end of the trajectory, which is the same thing only when labels are terminal, as they are in nearly every published temporal-recourse evaluation. Recourse feasibility is therefore a property of the benchmark configuration as much as of the method, and a validity number reported without its intervention-to-label distance is not interpretable. This is not a claim that recourse is impossible; it is a condition that every temporal recourse evaluation, including favourable ones,
needs to state.

\paragraph{Contributions.}
\begin{enumerate}
  \item \textbf{A model-vs-world audit instrument} for temporal counterfactual explanations: the proposal's own intervention is realised by the true mechanism, and the disagreement is decomposed additively into trajectory and outcome terms, each computed exactly against a known mechanism rather than estimated. Its construct validity is carried by do-complexity $D$, published beside every disagreement figure.
  \item \textbf{Intervention-to-outcome distance as a required reporting condition}, with the label site, not the trajectory end, established as the correct distance, using benchmark configurations that differ in exactly that field.
  \item \textbf{A benchmark suite with exact structural counterfactuals}: three additive-noise temporal SCM families shipping the true graph, the true mechanism, and oracle abduction--action--prediction counterfactuals, with the label functional exposed as a configurable axis rather than fixed by construction.
  \item \textbf{An audit of published temporal CF methods} against that instrument. \todo{one sentence, once results are final.}
\end{enumerate}

\paragraph{Scope.} Findings are established on the synthetic tier, which is the only tier possessing a ground-truth mechanism and therefore the only tier where a metric can be shown faithful to what it claims to measure. Real-signal tiers test whether the findings \emph{travel}; by construction they cannot validate the metrics, and we are explicit about not conflating the two. This distinction is the whole reason the instrument is trustworthy, and it is the same distinction the evaluation literature is currently missing.

\section{Related Work}

\paragraph{Counterfactual explanation for time series.}
Temporal counterfactual methods span instance substitution (Native Guide, \citealp{delaney2021nativeguide}), segment-level substitution (CoMTE, \citealp{ates2021comte}), evolutionary search (TSEvo, \citealp{hollig2022tsevo}), saliency-guided perturbation (M-CELS, \citealp{li2024mcels}), multi-objective optimisation (CONFETTI, \citealp{paredescetina2025confetti}), symbolic abstraction \citep{pludowski2025mascots}, and self-interpretable models that emit counterfactuals for their own predictions (CounTS, \citealp{yan2023counts}). Recent work has pushed on plausibility and temporal coherence of the edits themselves \citep{kostrzewa2026plausibility,zuin2026navigating,voets2026contex}.
What unifies the evaluation across this literature is that every reported quantity is model-relative: whether the classifier's decision changed, and how much of the input had to move for it to change. None of these papers is in a position to ask whether the proposed change would have worked, because their datasets do not come with a mechanism that could answer.

\paragraph{Benchmarks for temporal explanation.}
XTSC-Bench \citep{hollig2023xtscbench}, AMEE \citep{nguyen2024amee} and ExplainTS \citep{mozolewski2026explaints} standardise the comparison of temporal explainers, and the tabular CARLA library \citep{pawelczyk2021carla} to be distinguished from the algorithmic-recourse framework that shares its name \citep{karimi2021recourse} fixed the four-axis protocol the temporal work inherited. ExplainTS states the limitation directly: with real classification tasks it has no ground-truth relevance structure, so faithfulness must be assessed by proxy. Robustness and fragility critiques of explanation methods \citep{ghorbani2019fragile,yeh2019infidelity,artelt2021robustness} and of unjustified counterfactuals \citep{laugel2019dangers,poyiadzi2020face} sharpen the model-relative axes but do not add a world-relative one.

\paragraph{Causal temporal benchmarks.}
CausalTime \citep{cheng2024causaltime}, CausalDynamics \citep{herdeanu2025causaldynamics}, TimeGraph \citep{ferdous2025timegraph} and TCS \citep{gkorgkolis2025tcs} supply time series with known causal graphs, and CauKer \citep{xie2026cauker} and the causal-DAG prior of \citet{orourke2026causaldag} pair a ground-truth DAG with classification labels. All of them target \emph{structure recovery} or \emph{pretraining}: the object being scored is a discovered graph or a learned representation, never an independently produced post-hoc explanation. Their causal graphs are also typically atemporal over whole-trajectory nodes, rather than lagged over time-indexed variables, which is the structure an intervention at a specific timestep needs. DoFlow \citep{doflow2025} is the closest work on the counterfactual axis a flow-based model supporting observational, interventional and counterfactual queries over a known DAG but it is a generative forecaster in its own right, not a harness that scores third-party explainers against its ground truth. \citet{thumm2026interventional} make the relevant distinction explicit in their own appendix: their observational and interventional runs draw independent noise realisations, answering the \emph{interventional} query $P(X \mid \dop(\cdot))$ rather than the \emph{counterfactual} $P(X^{\mathrm{cf}} \mid X^{\mathrm{obs}}, \dop(\cdot))$, and they describe sharing the noise tensor exactly the step our oracle takes as a natural but strictly harder extension. Independently, they report the same attenuation of causal effects we characterise in Section~\ref{sec:audit}, arriving at it from a different stability mechanism and responding by restricting queries to short horizons rather than measuring the wall.

\paragraph{Causality used to constrain, versus causality used to derive.}
The causally-motivated counterfactual literature divides in a way that matters for what our benchmark can be. One branch uses a structural model to \textbf{penalise} implausible edits: the state of the art for temporal classification, \citet{bahri2025causalfeasibility}, adds a residual
$\lVert T'_t - f(\Pa(T'_t))\rVert$ to a Wachter-style objective over a free perturbation $T' = T + \delta$, no $\dop()$ operator, no intervention time, no abduction. That penalty is a \emph{local, one-step} check: it asks whether each slice is close to $f$ of its own parents, independently, and never propagates the intervention forward. A counterfactual can therefore be locally mechanism-consistent at every timestep while its intervention has no effect whatsoever on the outcome, and score perfectly. The second branch \textbf{derives} counterfactuals from a causal model, but a \emph{learned} one, in order to explain that same model's own predictions: CounTS \citep{yan2023counts} performs genuine abduction--action--prediction over the latent variables of its own variational model, and the static recourse literature \citep{karimi2021recourse} derives interventions from an assumed tabular SCM. Neither branch derives the counterfactual from a \emph{known} mechanism in order to score an \emph{independently produced} explanation against it. That conjunction ground-truth rather than learned mechanism, a graph lagged over time-indexed variables rather than atemporal over whole trajectories, and the derived counterfactual used as an evaluation oracle rather than as generative machinery is what this work adds.

\paragraph{Causal discovery.}
DYNOTEARS \citep{pamfil2020dynotears}, PCMCI \citep{runge2019pcmci} and the broader temporal causal-inference literature \citep{runge2023causal,lippe2022citris,song2024ctrlns} address the inverse problem: recovering structure from observations. In our synthetic tiers the graph is the generator's ground truth and discovery runs only as a diagnostic, so discovery error never enters a metric; we report it as a property of the benchmark, not of any explainer. Where a graph must be discovered rather than given, discovery accuracy becomes the ceiling on everything scored against it, and we say so explicitly rather than absorbing it silently into a method's score.

\section{The Model-vs-World Audit}
\label{sec:audit}


\todo{Populate. This section states the audit protocol and must cover, in
order:}
\begin{itemize}
  \item \todo{Setup and notation: temporal SCM, lagged graph
    $G \in \{0,1\}^{k\times k\times L}$, additive noise, the label functional,
    and the abduction-by-subtraction identity
    $\beps_t = \bx_t - f(\Pa(\bx_t))$ that everything downstream depends on.}
  \item \todo{The oracle structural counterfactual: abduction, action,
    prediction, with the noise tensor \emph{shared} between factual and
    counterfactual rollout (this is the step that makes the query
    counterfactual rather than interventional).}
  \item \todo{The admission gate (mechanism consistency), its two mutually
    exclusive semantics --- noiseless rollout versus Pearl delta-recursion ---
    and why a single counterfactual cannot satisfy both. State plainly that
    this is a pass/fail gate, not a ranking axis.}
  \item \todo{Extracting the intervention from a proposal, and do-complexity
    $D$ as the reading-invariant statistic; the two endpoint readings and why
    only the pair $(D, \Delta)$ is reportable.}
  \item \todo{The decomposition
    $\dototal = \dotraj + \dooutcome$, with the three outcome/trajectory
    combinations that define it, and the sign convention.}
  \item \todo{State the population the audit is computed over (flip-candidate
    instances) explicitly, so the reported disagreement is not read as a
    quantity averaged over a mixture. Probabilities of necessity and
    sufficiency are out of scope for this paper --- see Limitations.}
  \item \todo{The attenuation argument: spectral radius $<1$ implies geometric
    decay of an intervention's effect in $t_{\mathrm{label}} - t_0$, and the
    reporting rule that follows.}
\end{itemize}



\section{Submission of papers to NeurIPS 2026}


Please read the instructions below carefully and follow them faithfully.


\subsection{Style}


Papers to be submitted to NeurIPS 2026 must be prepared according to the
instructions presented here. Papers may only be up to {\bf nine} pages long, including figures. \textbf{Papers that exceed the page limit will not be reviewed (or in any other way considered) for presentation at the conference.}
Additional pages \emph{containing acknowledgments, references, checklist, and optional technical appendices} do not count as content pages. 


The margins in 2026 are the same as those in previous years.


Authors are required to use the NeurIPS \LaTeX{} style files obtainable at the NeurIPS website as indicated below. Please make sure you use the current files and not previous versions. Tweaking the style files may be grounds for desk rejection.


\subsection{Retrieval of style files}


The style files for NeurIPS and other conference information are available on the website at
\begin{center}
  \url{https://neurips.cc}.
\end{center}
% The file \verb+neurips_2026.pdf+ contains these instructions and illustrates the various formatting requirements your NeurIPS paper must satisfy.


The only supported style file for NeurIPS 2026 is \verb+neurips_2026.sty+, rewritten for \LaTeXe{}. \textbf{Previous style files for \LaTeX{} 2.09,  Microsoft Word, and RTF are no longer supported.}


The \LaTeX{} style file contains three optional arguments: 
\begin{itemize}
\item \verb+final+, which creates a camera-ready copy, 
\item \verb+preprint+, which creates a preprint for submission to, e.g., arXiv, \item \verb+nonatbib+, which will not load the \verb+natbib+ package for you in case of package clash.
\end{itemize}


\paragraph{Preprint option}
If you wish to post a preprint of your work online, e.g., on arXiv, using the NeurIPS style, please use the \verb+preprint+ option. This will create a nonanonymized version of your work with the text ``Preprint. Work in progress.'' in the footer. This version may be distributed as you see fit, as long as you do not say which conference it was submitted to. Please \textbf{do not} use the \verb+final+ option, which should \textbf{only} be used for papers accepted to NeurIPS.


At submission time, please omit the \verb+final+ and \verb+preprint+ options. This will anonymize your submission and add line numbers to aid review. Please do \emph{not} refer to these line numbers in your paper as they will be removed during generation of camera-ready copies.


The file \verb+neurips_2026.tex+ may be used as a ``shell'' for writing your paper. All you have to do is replace the author, title, abstract, and text of the paper with your own.


The formatting instructions contained in these style files are summarized in Sections \ref{gen_inst}, \ref{headings}, and \ref{others} below.


\section{General formatting instructions}
\label{gen_inst}


The text must be confined within a rectangle 5.5~inches (33~picas) wide and
9~inches (54~picas) long. The left margin is 1.5~inch (9~picas).  Use 10~point
type with a vertical spacing (leading) of 11~points.  Times New Roman is the
preferred typeface throughout, and will be selected for you by default.
Paragraphs are separated by \nicefrac{1}{2}~line space (5.5 points), with no
indentation.


The paper title should be 17~point, initial caps/lower case, bold, centered
between two horizontal rules. The top rule should be 4~points thick and the
bottom rule should be 1~point thick. Allow \nicefrac{1}{4}~inch space above and
below the title to rules. All pages should start at 1~inch (6~picas) from the
top of the page.


For the final version, authors' names are set in boldface, and each name is
centered above the corresponding address. The lead author's name is to be listed
first (left-most), and the co-authors' names (if different address) are set to
follow. If there is only one co-author, list both author and co-author side by
side.


Please pay special attention to the instructions in Section \ref{others}
regarding figures, tables, acknowledgments, and references.

\section{Headings: first level}
\label{headings}


All headings should be lower case (except for first word and proper nouns),
flush left, and bold.


First-level headings should be in 12-point type.


\subsection{Headings: second level}


Second-level headings should be in 10-point type.


\subsubsection{Headings: third level}


Third-level headings should be in 10-point type.


\paragraph{Paragraphs}


There is also a \verb+\paragraph+ command available, which sets the heading in
bold, flush left, and inline with the text, with the heading followed by 1\,em
of space.


\section{Citations, figures, tables, references}
\label{others}


These instructions apply to everyone.


\subsection{Citations within the text}


The \verb+natbib+ package will be loaded for you by default.  Citations may be
author/year or numeric, as long as you maintain internal consistency.  As to the
format of the references themselves, any style is acceptable as long as it is
used consistently.


The documentation for \verb+natbib+ may be found at
\begin{center}
  \url{http://mirrors.ctan.org/macros/latex/contrib/natbib/natnotes.pdf}
\end{center}
Of note is the command \verb+\citet+, which produces citations appropriate for
use in inline text.  For example,
\begin{verbatim}
   \citet{hasselmo} investigated\dots
\end{verbatim}
produces
\begin{quote}
  Hasselmo, et al.\ (1995) investigated\dots
\end{quote}


If you wish to load the \verb+natbib+ package with options, you may add the
following before loading the \verb+neurips_2026+ package:
\begin{verbatim}
   \PassOptionsToPackage{options}{natbib}
\end{verbatim}


If \verb+natbib+ clashes with another package you load, you can add the optional
argument \verb+nonatbib+ when loading the style file:
\begin{verbatim}
   \usepackage[nonatbib]{neurips_2026}
\end{verbatim}


As submission is double blind, refer to your own published work in the third
person. That is, use ``In the previous work of Jones et al.\ [4],'' not ``In our
previous work [4].'' If you cite your other papers that are not widely available
(e.g., a journal paper under review), use anonymous author names in the
citation, e.g., an author of the form ``A.\ Anonymous'' and include a copy of the anonymized paper in the supplementary material.


\subsection{Footnotes}


Footnotes should be used sparingly.  If you do require a footnote, indicate
footnotes with a number\footnote{Sample of the first footnote.} in the
text. Place the footnotes at the bottom of the page on which they appear.
Precede the footnote with a horizontal rule of 2~inches (12~picas).


Note that footnotes are properly typeset \emph{after} punctuation
marks.\footnote{As in this example.}


\subsection{Figures}


\begin{figure}
  \centering
  \fbox{\rule[-.5cm]{0cm}{4cm} \rule[-.5cm]{4cm}{0cm}}
  \caption{Sample figure caption. Explain what the figure shows and add a key take-away message to the caption.}
\end{figure}


All artwork must be neat, clean, and legible. Lines should be dark enough for
 reproduction purposes. The figure number and caption always appear after the
figure. Place one line space before the figure caption and one line space after
the figure. The figure caption should be lower case (except for the first word and proper nouns); figures are numbered consecutively.


You may use color figures.  However, it is best for the figure captions and the
paper body to be legible if the paper is printed in either black/white or in
color.


\subsection{Tables}


All tables must be centered, neat, clean, and legible.  The table number and
title always appear before the table.  See Table~\ref{sample-table}.


Place one line space before the table title, one line space after the
table title, and one line space after the table. The table title must
be lower case (except for the first word and proper nouns); tables are
numbered consecutively.


Note that publication-quality tables \emph{do not contain vertical rules}. We
strongly suggest the use of the \verb+booktabs+ package, which allows for
typesetting high-quality, professional tables:
\begin{center}
  \url{https://www.ctan.org/pkg/booktabs}
\end{center}
This package was used to typeset Table~\ref{sample-table}.


\begin{table}
  \caption{Sample table caption. Explain what the table shows and add a key take-away message to the caption.}
  \label{sample-table}
  \centering
  \begin{tabular}{lll}
    \toprule
    \multicolumn{2}{c}{Part}                   \\
    \cmidrule(r){1-2}
    Name     & Description     & Size ($\mu$m) \\
    \midrule
    Dendrite & Input terminal  & $\approx$100     \\
    Axon     & Output terminal & $\approx$10      \\
    Soma     & Cell body       & up to $10^6$  \\
    \bottomrule
  \end{tabular}
\end{table}

\subsection{Math}
Note that display math in bare TeX commands will not create correct line numbers for submission. Please use LaTeX (or AMSTeX) commands for unnumbered display math. (You really shouldn't be using \$\$ anyway; see \url{https://tex.stackexchange.com/questions/503/why-is-preferable-to} and \url{https://tex.stackexchange.com/questions/40492/what-are-the-differences-between-align-equation-and-displaymath} for more information.)

\subsection{Final instructions}

Do not change any aspects of the formatting parameters in the style files.  In
particular, do not modify the width or length of the rectangle the text should
fit into, and do not change font sizes. Please note that pages should be
numbered.


\section{Preparing PDF files}


Please prepare submission files with paper size ``US Letter,'' and not, for
example, ``A4.''


Fonts were the main cause of problems in the past years. Your PDF file must only
contain Type 1 or Embedded TrueType fonts. Here are a few instructions to
achieve this.


\begin{itemize}


\item You should directly generate PDF files using \verb+pdflatex+.


\item You can check which fonts a PDF files uses.  In Acrobat Reader, select the
  menu Files$>$Document Properties$>$Fonts and select Show All Fonts. You can
  also use the program \verb+pdffonts+ which comes with \verb+xpdf+ and is
  available out-of-the-box on most Linux machines.


\item \verb+xfig+ ``patterned'' shapes are implemented with bitmap fonts.  Use
  "solid" shapes instead.


\item The \verb+\bbold+ package almost always uses bitmap fonts.  You should use
  the equivalent AMS Fonts:
\begin{verbatim}
   \usepackage{amsfonts}
\end{verbatim}
followed by, e.g., \verb+\mathbb{R}+, \verb+\mathbb{N}+, or \verb+\mathbb{C}+
for $\mathbb{R}$, $\mathbb{N}$ or $\mathbb{C}$.  You can also use the following
workaround for reals, natural and complex:
\begin{verbatim}
   \newcommand{\RR}{I\!\!R} %real numbers
   \newcommand{\Nat}{I\!\!N} %natural numbers
   \newcommand{\CC}{I\!\!\!\!C} %complex numbers
\end{verbatim}
Note that \verb+amsfonts+ is automatically loaded by the \verb+amssymb+ package.


\end{itemize}


If your file contains type 3 fonts or non embedded TrueType fonts, we will ask
you to fix it.


\subsection{Margins in \LaTeX{}}


Most of the margin problems come from figures positioned by hand using
\verb+\special+ or other commands. We suggest using the command
\verb+\includegraphics+ from the \verb+graphicx+ package. Always specify the
figure width as a multiple of the line width as in the example below:
\begin{verbatim}
   \usepackage[pdftex]{graphicx} ...
   \includegraphics[width=0.8\linewidth]{myfile.pdf}
\end{verbatim}
See Section 4.4 in the graphics bundle documentation
(\url{http://mirrors.ctan.org/macros/latex/required/graphics/grfguide.pdf})


A number of width problems arise when \LaTeX{} cannot properly hyphenate a
line. Please give LaTeX hyphenation hints using the \verb+\-+ command when
necessary.

\begin{ack}
Use unnumbered first level headings for the acknowledgments. All acknowledgments
go at the end of the paper before the list of references. Moreover, you are required to declare
funding (financial activities supporting the submitted work) and competing interests (related financial activities outside the submitted work).
More information about this disclosure can be found at: \url{https://neurips.cc/Conferences/2026/PaperInformation/FundingDisclosure}.


Do {\bf not} include this section in the anonymized submission, only in the final paper. You can use the \texttt{ack} environment provided in the style file to automatically hide this section in the anonymized submission.
\end{ack}

\bibliographystyle{ieeetr}
\bibliography{references}
% \section*{References}


% References follow the acknowledgments in the camera-ready paper. Use unnumbered first-level heading for
% the references. Any choice of citation style is acceptable as long as you are
% consistent. It is permissible to reduce the font size to \verb+small+ (9 point)
% when listing the references.
% Note that the Reference section does not count towards the page limit.
% \medskip


% {
% \small


% [1] Alexander, J.A.\ \& Mozer, M.C.\ (1995) Template-based algorithms for
% connectionist rule extraction. In G.\ Tesauro, D.S.\ Touretzky and T.K.\ Leen
% (eds.), {\it Advances in Neural Information Processing Systems 7},
% pp.\ 609--616. Cambridge, MA: MIT Press.


% [2] Bower, J.M.\ \& Beeman, D.\ (1995) {\it The Book of GENESIS: Exploring
%   Realistic Neural Models with the GEneral NEural SImulation System.}  New York:
% TELOS/Springer--Verlag.


% [3] Hasselmo, M.E., Schnell, E.\ \& Barkai, E.\ (1995) Dynamics of learning and
% recall at excitatory recurrent synapses and cholinergic modulation in rat
% hippocampal region CA3. {\it Journal of Neuroscience} {\bf 15}(7):5249-5262.
% }


%%%%%%%%%%%%%%%%%%%%%%%%%%%%%%%%%%%%%%%%%%%%%%%%%%%%%%%%%%%%

\appendix

\section{Technical appendices and supplementary material}
Technical appendices with additional results, figures, graphs, and proofs may be submitted with the paper submission before the full submission deadline (see above). You can upload a ZIP file for videos or code, but do not upload a separate PDF file for the appendix. There is no page limit for the technical appendices.

Note: Think of the appendix as ``optional reading'' for reviewers. The paper must be able to stand alone without the appendix; for example, adding critical experiments that support the main claims to an appendix is inappropriate.

\end{document}